\documentclass[preprint]{revtex4-2}
\usepackage{amsmath}
\usepackage{graphicx}
\usepackage{comment}
\newcommand{\aap}{Astron. Ap.}
\newcommand{\apjs}{Astrophys. J. Suppl.}
\newcommand{\gca}{Geochim. Cosmochim. Acta}
\newcommand{\epsl}{Earth Planet. Sci. Lett.}

\begin{document}

\title{NRLEE Galactic Chemical Evolution\footnote{Supplemental material for the paper
{\it Constraints on the building blocks of the Solar System: Insights from nucleosynthetic isotope anomalies}, GCA (2021), authors: K. Bermingham,
B. S. Meyer, and K. Mezger, submitted}}
\author{Bradley S. Meyer}
\affiliation{Department of Physics and Astronomy, Clemson University, Clemson,
SC 29634, USA}
\author{Katherine R. Bermingham}
\affiliation{ Department of Earth and Planetary Sciences, Rutgers University, Piscataway, NJ 08854, USA}
\author{Klaus Mezger}
\affiliation{Institut f\"ur Geologie, Baltzerstrasse 1+3 3012 Bern, Schweiz}

\begin{abstract}
Readers may explore the results of this model in more detail via Jupyter
notebooks available at \url{http://webnucleo.org}.
\end{abstract}

\maketitle

\section{Introduction}\label{sec:intro}

\section{Galactic Chemical Model} \label{sec:initial}

\section{Inferred Anomalies} \label{sec:polytrope}

\section{Jupyter Notebooks} \label{sec:jupyter}

Readers who wish to study the above results in 
more detail are encouraged to explore the Jupyter notebooks accompanying
the Webnucleo projects used to construct the simple NRLEE model presented
here.  The notebooks allow a user to download and graph data interactively.
Input parameters can be changed, and figures and tables saved.  The
notebooks run on Google Colab (\url{https://colab.research.google.com)} or
Binder (\url{https://mybinder.org}).  They can also be downloaded and run
on a local computer.

The notebooks relevant for the model presented here are:
\begin{itemize}
\item Mainline-nucleosynthesis: \url{https://github.com/mbradle/mainline-nucleosynthesis}
\item simple\_s\_process: \url{https://github.com/mbradle/simple_s_process}
\item NRLEE-Nucleosynthesis: \url{https://github.com/mbradle/NRLEE-Nucleosynthesis}
\end{itemize}
All of these notebooks are part of the Webnucleo Project
(\url{http://webnucleo.org}).

\section{Conclusion}

\bibliography{clemson}

\begin{figure}
\includegraphics{figures/gas_frac}
\caption{The dust and gas fraction as a function of Galactic time in the model.}
\label{fig:gas_frac}
\end{figure}

\begin{figure}
\includegraphics{figures/dust_mass}
\caption{The mass in various dust components.}
\label{fig:dust}
\end{figure}

\begin{figure}
\includegraphics{figures/x}
\caption{The mass fractions of species in the molecular cloud phase of the model
as a function of Galactic time.}
\label{fig:x}
\end{figure}

\begin{figure}
\includegraphics{figures/ca}
\caption{The anomailes in Ca.}
\label{fig:ca}
\end{figure}

\begin{figure}
\includegraphics{figures/ti}
\caption{The anomailes in Ti.}
\label{fig:ti}
\end{figure}

\begin{figure}
\includegraphics{figures/cr}
\caption{The anomailes in Cr.}
\label{fig:cr}
\end{figure}

\begin{figure}
\includegraphics{figures/fe}
\caption{The anomailes in Fe.}
\label{fig:fe}
\end{figure}

\begin{figure}
\includegraphics{figures/ni}
\caption{The anomailes in Ni.}
\label{fig:ni}
\end{figure}

\begin{figure}
\includegraphics{figures/ca48_ti50}
\caption{The correlated anomalies in $^{48}$Ca and $^{50}$Ti.}
\label{fig:ca48_ti50}
\end{figure}

\begin{figure}
\includegraphics{figures/ti46_ca48}
\caption{The correlated anomalies in $^{46}$Ti and $^{48}$Ca.}
\label{fig:ti46_ca48}
\end{figure}

\begin{figure}
\includegraphics{figures/ti46_ti50}
\caption{The correlated anomalies in $^{46}$Ti and $^{50}$Ti.}
\label{fig:ti46_ti50}
\end{figure}

\begin{figure}
\includegraphics{figures/cr54_ti50}
\caption{The correlated anomalies in $^{54}$Cr and $^{50}$Ti.}
\label{fig:cr54_ti50}
\end{figure}

%\begin{figure}
%\includegraphics{figures/ni64_cr54}
%\caption{The correlated anomalies in $^{54}$Cr and $^{64}$Ni.}
%\label{fig:cr54_ni64}
%\end{figure}

\end{document}
